\SourceLeaf{531}
\begin{center}
\begin{tabular}{rrr@{\hspace{2em}}rrr}
Thousand & \textsc{Gauss} & True Value & Thousand & \textsc{Gauss} & True Value \\
$20$ & $102$ & $104$ & $546$ & $68$ & $69$ \\
$159$ & $87$ & $77$ & $601$ & $75$ & $76$ \\
$199$ & $96$ & $86$ & $625$ & $68$ & $78$ \\
$206$ & $85$ & $83$ & $668$ & $73$ & $74$ \\
$245$ & $78$ & $88$ & $675$ & $69$ & $73$ \\
$289$ & $85$ & $77$ & $784$ & $74$ & $75$ \\
$290$ & $84$ & $85$ & $800$ & $81$ & $71$ \\
$334$ & $80$ & $81$ & $879$ & $68$ & $78$ \\
$352$ & $80$ & $81$ & $985$ & $74$ & $70$ \\
$501$ & $78$ & $79$ &  &  &  \\
\end{tabular}
\end{center}
\par The formula of \textsc{Encke} mentioned in the letter from \textsc{Gauss} to \textsc{Encke} seems to be the following
\[
\frac{n}{\log n}\overset{2\log n}{\surd}10
\]
which \textsc{Encke} gives in a letter to \textsc{Gauss} of 4 December 1849.
\par The \textit{Table of the Number of Classes of Binary Quadratic Forms} gives the number of genera and classes and the index of irregularity for negative determinants in hundred groups 1 to 30, 43, 51, 61, 62, 63, 91 to 100, 117 to 120; also, in a separate compilation, for those of the 1st, 3rd and 10th thousand groups, for the first 800 of the forms $-(15n+7)$ and $-(15n+13)$, and for some very large determinants; and for positive determinants in the 1st, 2nd, 3rd, 9th and 10th hundred groups and some others. The manuscript consists of separate slips with tables arranged in various ways. For instance, the older ones use the word Ordo in place of Genus, as in the individual hundred groups except the 9th and 10th for positive determinants, and also in some preliminary compilations in thousand groups.
To make the arrangement easier to survey, the notation of the Disqu. Arithm. has been used throughout here. For lack of space, the greatest and least quotients obtained by dividing the number of classes by the determinant, and the numbers of determinants for which the quotient falls within certain bounds, are not placed below the individual hundred groups, but collected at the end of the table for negative determinants. Some surviving notes seem to indicate that \textsc{Gauss} first calculated the classes for determinants belonging to the same hundred group and the same remainder on division by 15. The determinants in these divisions were then arranged by the number of genera and classes, and finally all those in the same hundred group were ordered in the way reproduced here.
Several later corrections have been inserted into the tables of individual hundred groups, but not into the compilations in thousand groups. Only the two tables with determinants of the forms $-(15n+7)$ and $-(15n+13)$ contain dates, respectively ‘\textit{Expl. In. Febr. 1801}’ and ‘\textit{Expl. 27 Febr. 1807.}’
\par Among the errors I have noticed in these tables are the following. To make them easier to check, I give the period numbers of the fundamental classes, for example $4\pdot\ 4\pdot\ 2$ for the determinant
